\chapter{Reproducing the results}
\label{app:eval}

This appendix explains how to run the laboratory and regenerate every result reported in Chapter~\ref{cha:eval}. Commands are run from the repository root.

\section{Environment}
\label{sec:run-env}

The results were produced with Node.js~22, Python~3, PyTorch~2.14.1 (recorded in \path{adversarial_metrics.json}) and scikit-learn; the phishing results were also reproduced independently with scikit-learn~1.9.1, with byte-identical output. The adversarial results were produced on a CPU. The CIFAR-10 test set must be present under \path{eval/adversarial/results/cifar10/}; the model weights are downloaded by \texttt{torch.hub} on first use.

\section{Running the laboratory}
\label{sec:run-web}

\begin{minted}{text}
cp .env.example .env.local
# set GEMINI_API_KEY, LAB_BASIC_USER and LAB_BASIC_PASS
# optional: GEMINI_MODEL, EVAL_API_URL, EVAL_INTERNAL_TOKEN
npm install
npm test
python3 eval/adversarial/service.py &   # attack service on port 8001
npm run dev                              # http://localhost:3000
\end{minted}

For a production-style run, \texttt{npm run build} followed by \path{deploy/start-live-lab.sh} starts both processes and refuses to start without credentials. The container images are built from \texttt{Dockerfile} (web server only) and \path{deploy/Dockerfile.all-in-one} (web server and attack service).

\section{Phishing benchmark}
\label{sec:run-phish}

\begin{minted}{text}
python3 -m pip install -r eval/requirements.txt
python3 eval/phishing/run_benchmark.py
\end{minted}

The script regenerates \path{eval/phishing/corpus.jsonl} from code and writes \path{phishing_metrics.json} and \texttt{predictions.csv} under \path{eval/phishing/results/}. It does not overwrite \path{MANUAL_REVIEW.csv}.

\section{Adversarial batch}
\label{sec:run-adv}

\begin{minted}{text}
python3 eval/adversarial/run_attacks.py \
  2>&1 | tee eval/adversarial/results/run200_full.txt
\end{minted}

The script writes \path{eval/adversarial/results/adversarial_metrics.json}, including the commit and PyTorch version, and stops with an error if any perturbation exceeds its budget. The log of the reported run is \path{run200_full.txt}.

\section{Live attack calls}
\label{sec:run-live}

With the attack service running, the recorded live calls are reproduced with:

\begin{minted}{text}
curl -s -X POST localhost:8001/attack -H 'Content-Type: application/json' \
  -d '{"originalClass":"airplane","noiseLevel":0.05,"method":"fgsm"}'
curl -s -X POST localhost:8001/attack -H 'Content-Type: application/json' \
  -d '{"originalClass":"airplane","noiseLevel":0.05,"method":"pgd"}'
curl -s -X POST localhost:8001/attack -H 'Content-Type: application/json' \
  -d '{"originalClass":"airplane","targetClass":"bird","targeted":true,
       "noiseLevel":0.05,"method":"fgsm"}'
\end{minted}
